% ===========================================================================
%  Chapitre 1 — Équations et inéquations du second degré
%  PE 2026, composante ALGÈBRE
% ===========================================================================
\bmtchapitre{Équations et inéquations du second degré}
  {Résoudre dans $\R$ une équation et une inéquation du second degré,
   factoriser un trinôme, étudier son signe, et ramener au second degré une
   équation bicarrée, une équation du troisième degré ou une équation
   quotient.}
  {Algèbre \textbullet\ environ 6 heures}

Ce chapitre est le plus court du livre et c'est celui dont vous vous servirez
le plus. Non pas parce qu'on vous demandera souvent, en soi, de résoudre
$2x^{2}-3x+1 = 0$ : parce que ce calcul-là surgira à l'intérieur de presque
tout le reste. Le signe d'une dérivée se lit sur un trinôme. L'ensemble de
définition d'une fonction rationnelle dépend des racines de son dénominateur.
Une droite d'ajustement se cherche en résolvant un système. Le second degré
est l'outil, pas le sujet.

Vous l'avez déjà rencontré en seconde et en première. Ce qui change cette
année, c'est la vitesse : le discriminant, les racines, le tableau de signes
doivent sortir sans hésitation, en moins d'une minute, pour que la tête reste
disponible pour la question qui suit.

% ---------------------------------------------------------------------------
\begin{revision}
Rien de nouveau ici : ce sont les outils de seconde et de première dont le
chapitre va se servir.

\begin{enumerate}[label=\textbf{\arabic*.}, leftmargin=*, itemsep=0.35em]
  \item Développer $(x-3)^{2}$ et $(2x+5)(2x-5)$.
  \item Factoriser $x^{2}-9$ puis $x^{2}+6x+9$.
  \item Résoudre $3x-7 = 0$, puis $-2x+8 \geq 0$.
  \item Donner le signe de $(x-1)(x+4)$ selon les valeurs de $x$.
  \item Pour quelles valeurs de $x$ le quotient $\dfrac{1}{x-2}$
        existe-t-il ?
  \item Un produit de deux facteurs est nul. Que peut-on en conclure ?
\end{enumerate}
\end{revision}

\begin{automatismes}
Sans calculatrice, sans brouillon.

\begin{enumerate}[label=\textbf{\alph*)}, leftmargin=*, itemsep=0.25em]
  \item $(-3)^{2}$
  \item $\sqrt{49}$
  \item $\sqrt{0{,}25}$
  \item $b^{2}-4ac$ pour $a=1$, $b=5$, $c=6$
  \item Les solutions de $x^{2} = 16$
  \item Les solutions de $x(x-7) = 0$
  \item Le signe de $x^{2}+1$
  \item $\dfrac{-6+2}{2\times 2}$
\end{enumerate}
\end{automatismes}

\begin{corrige}[automatismes]
a) $9$ \quad b) $7$ \quad c) $0{,}5$ \quad d) $1$ \quad e) $-4$ et $4$ \quad
f) $0$ et $7$ \quad g) strictement positif pour tout réel \quad
h) $-1$.
\end{corrige}

% ===========================================================================
\section{Le trinôme et son discriminant}

\begin{definition}[trinôme du second degré]
On appelle \textbf{trinôme du second degré} toute expression de la forme
\[
  P(x) = ax^{2}+bx+c, \qquad a,b,c \in \R, \quad a \neq 0 .
\]
La condition $a \neq 0$ n'est pas un détail : si $a$ était nul, l'expression
serait du premier degré et tout ce chapitre serait sans objet.
\end{definition}

\begin{propriete}[forme canonique]
Pour tout réel $x$,
\[
  ax^{2}+bx+c
  = a\left[\left(x+\frac{b}{2a}\right)^{\!2}-\frac{b^{2}-4ac}{4a^{2}}\right].
\]
\end{propriete}

\begin{demonstration}
On met $a$ en facteur, puis on complète le carré :
\[
  ax^{2}+bx+c = a\left(x^{2}+\frac{b}{a}x\right)+c
  = a\left[\left(x+\frac{b}{2a}\right)^{\!2}-\frac{b^{2}}{4a^{2}}\right]+c .
\]
En rentrant $c$ dans le crochet, $c = a\cdot\dfrac{4ac}{4a^{2}}$, ce qui donne
le résultat annoncé.
\end{demonstration}

\begin{definition}[discriminant]
Le nombre $\Delta = b^{2}-4ac$ s'appelle le \textbf{discriminant} du trinôme.
Le mot vient de \emph{discriminer}, distinguer : c'est lui, et lui seul, qui
distingue les trois cas possibles.
\end{definition}

\begin{theoreme}[résolution de $ax^{2}+bx+c = 0$]
Soit $\Delta = b^{2}-4ac$.
\begin{itemize}[leftmargin=1.3em, itemsep=0.2em]
  \item Si $\Delta > 0$, l'équation admet \textbf{deux} solutions distinctes
        \[
          x_{1} = \frac{-b-\sqrt{\Delta}}{2a}, \qquad
          x_{2} = \frac{-b+\sqrt{\Delta}}{2a} .
        \]
  \item Si $\Delta = 0$, l'équation admet \textbf{une} solution double
        $x_{0} = -\dfrac{b}{2a}$.
  \item Si $\Delta < 0$, l'équation n'admet \textbf{aucune} solution réelle.
\end{itemize}
\end{theoreme}

\begin{demonstration}
D'après la forme canonique, l'équation équivaut à
$\left(x+\frac{b}{2a}\right)^{2} = \dfrac{\Delta}{4a^{2}}$. Un carré est
toujours positif ou nul : si $\Delta<0$, aucun réel ne convient. Si
$\Delta = 0$, la seule possibilité est $x+\frac{b}{2a} = 0$. Si $\Delta>0$,
on obtient $x+\frac{b}{2a} = \pm\dfrac{\sqrt\Delta}{2a}$, d'où les deux
racines.
\end{demonstration}

\begin{visualisation}[les trois cas, vus d'en haut]
\begin{center}
\begin{tikzpicture}
\begin{axis}[
  width = 13.2cm, height = 5.6cm,
  axis lines = middle, xlabel = {$x$},
  xmin = -3.4, xmax = 3.4, ymin = -2.2, ymax = 4.2,
  xtick = \empty, ytick = \empty,
  axis line style = {bmtGris}, clip = true,
]
  \addplot[bmtLaterite, line width=1.3pt, domain=-3.2:-0.15, samples=80]
    {(x+2.4)*(x+0.8)};
  \addplot[bmtIndigo, line width=1.3pt, domain=-1.4:1.4, samples=80]
    {(x)^2};
  \addplot[bmtVetiver, line width=1.3pt, domain=1.05:3.3, samples=80]
    {(x-2.2)^2+1.1};
  \node[bmtLaterite, font=\footnotesize, anchor=south]
    at (axis cs:-1.6,2.3) {$\Delta > 0$};
  \node[bmtIndigo, font=\footnotesize, anchor=south]
    at (axis cs:0,2.3) {$\Delta = 0$};
  \node[bmtVetiver, font=\footnotesize, anchor=south]
    at (axis cs:2.2,3.0) {$\Delta < 0$};
  \filldraw[bmtLaterite] (axis cs:-2.4,0) circle (1.7pt);
  \filldraw[bmtLaterite] (axis cs:-0.8,0) circle (1.7pt);
  \filldraw[bmtIndigo] (axis cs:0,0) circle (1.7pt);
\end{axis}
\end{tikzpicture}
\end{center}

Résoudre $ax^{2}+bx+c = 0$, c'est chercher où la parabole rencontre l'axe des
abscisses. À gauche, elle le traverse en deux points : deux racines. Au
milieu, elle le touche sans le traverser : une racine double. À droite, elle
reste au-dessus : aucune racine réelle. Le discriminant est le calcul qui dit
lequel des trois dessins vous avez sous les yeux, sans avoir à le tracer.
\end{visualisation}

\begin{exemple}[les trois cas, déroulés]
\textbf{a)} $2x^{2}-3x-2 = 0$. Ici $a=2$, $b=-3$, $c=-2$, donc
$\Delta = 9+16 = 25 > 0$ et $\sqrt\Delta = 5$ :
\[
  x_{1} = \frac{3-5}{4} = -\frac12, \qquad x_{2} = \frac{3+5}{4} = 2 .
\]

\textbf{b)} $x^{2}-6x+9 = 0$. Ici $\Delta = 36-36 = 0$, donc une racine
double $x_{0} = \dfrac{6}{2} = 3$.

\textbf{c)} $x^{2}+x+1 = 0$. Ici $\Delta = 1-4 = -3 < 0$ : pas de solution
réelle.
\end{exemple}

\begin{methode}{résoudre une équation du second degré}
\begin{enumerate}[label=\textbf{\arabic*.}, leftmargin=*, itemsep=0.15em]
  \item \emph{Mettre en forme.} Tout ramener dans le membre de gauche, de
        sorte que l'équation s'écrive $ax^{2}+bx+c = 0$.
  \item \emph{Regarder avant de calculer.} Si $c = 0$, factoriser par $x$.
        Si $b = 0$, isoler $x^{2}$. Dans ces deux cas, le discriminant est
        inutile.
  \item \emph{Identifier} $a$, $b$, $c$ — avec leur signe.
  \item \emph{Calculer} $\Delta = b^{2}-4ac$, conclure selon son signe.
  \item \emph{Vérifier} : la somme des racines doit valoir $-\frac ba$
        (voir la section suivante). Ce contrôle prend cinq secondes.
\end{enumerate}
\end{methode}

\begin{erreurclassique}
$\Delta = b^{2}-4ac$, et non $b^{2}-4ac$ « où le signe moins n'est pas à
prendre en compte ». Pour $x^{2}-3x-4 = 0$, on a $c = -4$, donc
$-4ac = -4\times1\times(-4) = +16$ et $\Delta = 9+16 = 25$. L'erreur la plus
fréquente sur les copies consiste à écrire $\Delta = 9-16 = -7$ et à conclure
qu'il n'y a pas de solution, alors qu'il y en a deux : $-1$ et $4$.
\end{erreurclassique}

\begin{entrainement}[équations du second degré]
\exo[\niveau{1}] $x^{2}-5x+6 = 0$

\exo[\niveau{1}] $x^{2}+4x+4 = 0$

\exo[\niveau{1}] $3x^{2}-12 = 0$

\exo[\niveau{1}] $2x^{2}+7x = 0$

\exo[\niveau{2}] $-x^{2}+2x+8 = 0$

\exo[\niveau{2}] $\dfrac{1}{2}x^{2}-x-4 = 0$

\exo[\niveau{2}] $x^{2}+x+3 = 0$

\exo[\niveau{2}] $(x-1)^{2} = 2x+7$

\exo[\niveau{3}] Déterminer le réel $m$ pour que l'équation
$x^{2}-4x+m = 0$ admette une racine double, puis donner cette racine.
\end{entrainement}

\begin{corrige}[équations du second degré]
\textbf{1.} $\Delta = 1$ ; $x_{1} = 2$, $x_{2} = 3$.
\textbf{2.} $\Delta = 0$ ; $x_{0} = -2$.
\textbf{3.} $x^{2} = 4$ : $x = -2$ ou $x = 2$ (pas de discriminant à
calculer).
\textbf{4.} $x(2x+7) = 0$ : $x = 0$ ou $x = -\frac72$.
\textbf{5.} $\Delta = 4+32 = 36$ ; $x_{1} = \frac{-2-6}{-2} = 4$,
$x_{2} = \frac{-2+6}{-2} = -2$.
\textbf{6.} $\Delta = 1+8 = 9$ ; $x_{1} = \frac{1-3}{1} = -2$,
$x_{2} = \frac{1+3}{1} = 4$.
\textbf{7.} $\Delta = 1-12 = -11 < 0$ : aucune solution réelle.
\textbf{8.} $x^{2}-2x+1 = 2x+7$ donne $x^{2}-4x-6 = 0$, $\Delta = 40$,
$x = 2\pm\sqrt{10}$.
\textbf{9.} Racine double $\iff \Delta = 0 \iff 16-4m = 0 \iff m = 4$ ; la
racine est alors $x_{0} = 2$.
\end{corrige}

% ===========================================================================
\section{Somme, produit et factorisation}

\begin{propriete}[somme et produit des racines]
Si $\Delta \geq 0$ et si $x_{1}$, $x_{2}$ désignent les racines (confondues
lorsque $\Delta = 0$), alors
\[
  S = x_{1}+x_{2} = -\frac{b}{a}, \qquad P = x_{1}\,x_{2} = \frac{c}{a}.
\]
\end{propriete}

\begin{demonstration}
En additionnant les deux expressions des racines, les $\pm\sqrt\Delta$
s'éliminent : $x_{1}+x_{2} = \dfrac{-2b}{2a} = -\dfrac ba$. En les
multipliant, on reconnaît une identité remarquable :
\[
  x_{1}x_{2} = \frac{(-b)^{2}-\left(\sqrt\Delta\right)^{2}}{4a^{2}}
  = \frac{b^{2}-(b^{2}-4ac)}{4a^{2}} = \frac{c}{a}. \qedhere
\]
\end{demonstration}

\begin{theoreme}[factorisation d'un trinôme]
Si $\Delta > 0$, alors $ax^{2}+bx+c = a\,(x-x_{1})(x-x_{2})$.
Si $\Delta = 0$, alors $ax^{2}+bx+c = a\,(x-x_{0})^{2}$.
Si $\Delta < 0$, le trinôme ne se factorise pas dans $\R$.
\end{theoreme}

\begin{remarque}
Cette factorisation est la raison pour laquelle ce chapitre revient sans
arrêt. Chaque fois qu'un quotient contient un trinôme, c'est elle qui permet
de simplifier ; chaque fois qu'un signe est demandé, c'est elle qui le donne.
\end{remarque}

\begin{exemple}[deux nombres dont on connaît la somme et le produit]
Cherchons deux nombres dont la somme vaut $7$ et le produit $12$. Ce sont les
racines de $x^{2}-7x+12 = 0$, car pour ce trinôme $S = 7$ et $P = 12$. Son
discriminant vaut $49-48 = 1$, d'où $x_{1} = 3$ et $x_{2} = 4$. On vérifie :
$3+4 = 7$ et $3\times4 = 12$.
\end{exemple}

\begin{entrainement}[somme, produit, factorisation]
\exo[\niveau{1}] Factoriser $x^{2}-x-6$.

\exo[\niveau{1}] Factoriser $2x^{2}-8x+8$.

\exo[\niveau{2}] Sans calculer les racines, donner la somme et le produit des
racines de $5x^{2}-3x-2 = 0$. En déduire une racine évidente, puis l'autre.

\exo[\niveau{2}] Trouver deux nombres dont la somme vaut $-1$ et le produit
$-20$.

\exo[\niveau{2}] Simplifier $\dfrac{x^{2}-4}{x^{2}-5x+6}$ pour $x \neq 2$ et
$x \neq 3$.

\exo[\niveau{3}] Le trinôme $x^{2}+bx+9$ admet une racine double. Déterminer
les valeurs possibles de $b$.
\end{entrainement}

\begin{corrige}[somme, produit, factorisation]
\textbf{1.} $\Delta = 25$, racines $-2$ et $3$ : $(x+2)(x-3)$.
\textbf{2.} $\Delta = 0$, racine double $2$ : $2(x-2)^{2}$.
\textbf{3.} $S = \frac35$, $P = -\frac25$. On voit que $x = 1$ convient
($5-3-2=0$) ; l'autre racine vaut alors $P/1 = -\frac25$.
\textbf{4.} Racines de $x^{2}+x-20 = 0$ : $\Delta = 81$, soit $-5$ et $4$.
\textbf{5.} $\dfrac{(x-2)(x+2)}{(x-2)(x-3)} = \dfrac{x+2}{x-3}$.
\textbf{6.} $\Delta = b^{2}-36 = 0$ donne $b = -6$ ou $b = 6$.
\end{corrige}

% ===========================================================================
\section{Signe du trinôme et inéquations}

\begin{theoreme}[signe d'un trinôme]
Un trinôme est \textbf{du signe de $a$ partout}, sauf entre ses racines,
lorsqu'elles existent, où il est du signe contraire.
\end{theoreme}

\begin{demonstration}
Si $\Delta<0$, la forme canonique montre que le crochet est strictement
positif : le trinôme a le signe de $a$ pour tout $x$. Si $\Delta \geq 0$, on
écrit $a(x-x_{1})(x-x_{2})$ : le produit $(x-x_{1})(x-x_{2})$ est négatif
exactement lorsque $x$ est strictement compris entre les deux racines,
positif ailleurs.
\end{demonstration}

\begin{visualisation}[le tableau de signes que l'on écrira toute l'année]
\begin{center}
\begin{tikzpicture}[scale=1]
  \draw[bmtGris, line width=0.5pt] (0,0) rectangle (11.6,1.9);
  \draw[bmtGris, line width=0.5pt] (0,1.25) -- (11.6,1.25);
  \draw[bmtGris, line width=0.5pt] (2.1,0) -- (2.1,1.9);
  \node[font=\footnotesize] at (1.05,1.575) {$x$};
  \node[font=\footnotesize] at (1.05,0.62) {$ax^{2}+bx+c$};
  \node[font=\footnotesize] at (2.5,1.575) {$-\infty$};
  \node[font=\footnotesize] at (5.3,1.575) {$x_{1}$};
  \node[font=\footnotesize] at (8.3,1.575) {$x_{2}$};
  \node[font=\footnotesize] at (11.2,1.575) {$+\infty$};
  \draw[bmtGris, line width=0.5pt] (5.3,0) -- (5.3,1.25);
  \draw[bmtGris, line width=0.5pt] (8.3,0) -- (8.3,1.25);
  \node[bmtIndigo, font=\footnotesize] at (3.9,0.62) {signe de $a$};
  \node[bmtLaterite, font=\footnotesize] at (6.8,0.62) {signe de $-a$};
  \node[bmtIndigo, font=\footnotesize] at (9.75,0.62) {signe de $a$};
  \node[font=\footnotesize] at (5.3,0.62) {$0$};
  \node[font=\footnotesize] at (8.3,0.62) {$0$};
\end{tikzpicture}
\end{center}

Trois colonnes, deux zéros, et le signe qui alterne : voilà tout ce qu'il y a
à retenir. Lorsque $\Delta = 0$, les deux zéros se confondent et le signe
n'alterne pas — le trinôme s'annule en un point sans changer de signe.
Lorsque $\Delta < 0$, il n'y a ni zéro ni changement : une seule case, du
signe de $a$.
\end{visualisation}

\begin{methode}{résoudre une inéquation du second degré}
\begin{enumerate}[label=\textbf{\arabic*.}, leftmargin=*, itemsep=0.15em]
  \item Tout ramener à gauche : $ax^{2}+bx+c \geq 0$, ou $\leq 0$, etc.
  \item Chercher les racines du trinôme (discriminant).
  \item Dresser le tableau de signes.
  \item Lire l'ensemble solution, en soignant les crochets : fermés si
        l'inégalité est large, ouverts si elle est stricte.
\end{enumerate}
\end{methode}

\begin{exemple}[une inéquation, du calcul à l'écriture de la solution]
Résolvons $-x^{2}+3x+4 \geq 0$.

Le discriminant vaut $9+16 = 25$, les racines sont
$x_{1} = \frac{-3-5}{-2} = 4$ et $x_{2} = \frac{-3+5}{-2} = -1$. Ici
$a = -1 < 0$ : le trinôme est négatif à l'extérieur de $\intervalle{-1}{4}$
et positif entre les racines. L'inégalité étant large, les bornes sont
incluses :
\[
  \mathcal{S} = \intervalle{-1}{4}.
\]
\end{exemple}

\begin{erreurclassique}
Ordonner les racines avant de remplir le tableau. Dans l'exemple ci-dessus,
la formule donne d'abord $4$ puis $-1$ ; si l'on recopie ces deux nombres
dans cet ordre en tête du tableau, tous les signes sont faux. On écrit
toujours la plus petite racine à gauche.
\end{erreurclassique}

\begin{entrainement}[inéquations]
\exo[\niveau{1}] $x^{2}-4 > 0$

\exo[\niveau{1}] $x^{2}-5x+6 \leq 0$

\exo[\niveau{2}] $-2x^{2}+x+1 \geq 0$

\exo[\niveau{2}] $x^{2}+x+1 > 0$

\exo[\niveau{2}] $3x^{2}-6x \leq 0$

\exo[\niveau{3}] $\dfrac{x^{2}-1}{x-3} \leq 0$ \quad (attention à la valeur
interdite)

\exo[\niveau{3}] Pour quelles valeurs de $x$ le nombre $\sqrt{x^{2}-3x+2}$
est-il défini ?
\end{entrainement}

\begin{corrige}[inéquations]
\textbf{1.} Racines $-2$ et $2$, $a>0$ :
$\mathcal S = \intervalleo{-\infty}{-2}\cup\intervalleo{2}{+\infty}$.
\textbf{2.} Racines $2$ et $3$, $a>0$ : $\mathcal S = \intervalle{2}{3}$.
\textbf{3.} $\Delta = 1+8 = 9$, racines $-\frac12$ et $1$, $a<0$ :
$\mathcal S = \intervalle{-\frac12}{1}$.
\textbf{4.} $\Delta = -3<0$ et $a>0$ : le trinôme est toujours strictement
positif, $\mathcal S = \R$.
\textbf{5.} $3x(x-2) \leq 0$ : $\mathcal S = \intervalle{0}{2}$.
\textbf{6.} Valeur interdite : $3$. On dresse un tableau à trois lignes. Le
numérateur $x^{2}-1$ est positif à l'extérieur de $\intervalle{-1}{1}$,
négatif à l'intérieur ; le dénominateur $x-3$ est négatif avant $3$, positif
après. Le quotient est donc négatif sur $\intervalleo{-\infty}{-1}$, positif
sur $\intervalleo{-1}{1}$, négatif sur $\intervalleo{1}{3}$, positif sur
$\intervalleo{3}{+\infty}$. Les zéros $-1$ et $1$ sont acceptés,
$3$ est exclu :
$\mathcal S = \intervalleof{-\infty}{-1}\cup\intervallefo{1}{3}$.
\textbf{7.} La racine carrée exige $x^{2}-3x+2 \geq 0$, de racines $1$ et
$2$ avec $a>0$ :
$\mathcal S = \intervalleof{-\infty}{1}\cup\intervallefo{2}{+\infty}$.
\end{corrige}

\begin{qcm}
\exo Le discriminant de $-3x^{2}+2x-1$ vaut :
\textbf{a)} $16$ \quad \textbf{b)} $-8$ \quad \textbf{c)} $8$ \quad
\textbf{d)} $-16$ \reponseqcm{b}

\exo L'équation $x^{2}+5 = 0$ admet dans $\R$ :
\textbf{a)} deux solutions \quad \textbf{b)} une solution double \quad
\textbf{c)} aucune solution \quad \textbf{d)} une infinité \reponseqcm{c}

\exo Les racines de $x^{2}-2x-15$ sont $-3$ et $5$. La forme factorisée est :
\textbf{a)} $(x-3)(x+5)$ \quad \textbf{b)} $(x+3)(x-5)$ \quad
\textbf{c)} $(x-3)(x-5)$ \quad \textbf{d)} $-(x+3)(x-5)$ \reponseqcm{b}

\exo L'ensemble des solutions de $x^{2} \leq 9$ est :
\textbf{a)} $\intervalleo{-\infty}{3}$ \quad \textbf{b)} $\intervalle{-3}{3}$
\quad \textbf{c)} $\intervallefo{3}{+\infty}$ \quad
\textbf{d)} $\intervalleof{-\infty}{-3}\cup\intervallefo{3}{+\infty}$
\reponseqcm{b}

\exo Si $S = 6$ et $P = 8$ sont la somme et le produit des racines, le
trinôme unitaire correspondant est :
\textbf{a)} $x^{2}+6x+8$ \quad \textbf{b)} $x^{2}-6x-8$ \quad
\textbf{c)} $x^{2}-6x+8$ \quad \textbf{d)} $x^{2}+6x-8$ \reponseqcm{c}
\end{qcm}

% ===========================================================================
\section{Les équations qui se ramènent au second degré}

Trois familles reviennent régulièrement. Dans les trois cas, la démarche est
la même : \emph{ne pas} chercher à résoudre directement, mais transformer
l'équation jusqu'à faire apparaître un trinôme.

\subsection{L'équation bicarrée}

\begin{methode}{résoudre $ax^{4}+bx^{2}+c = 0$}
Poser $X = x^{2}$, avec la contrainte $X \geq 0$. Résoudre $aX^{2}+bX+c = 0$,
puis revenir : chaque solution $X$ strictement positive donne deux valeurs de
$x$, à savoir $\pm\sqrt X$ ; la solution $X = 0$ en donne une seule ; une
solution négative n'en donne aucune.
\end{methode}

\begin{exemple}[une bicarrée complète]
$x^{4}-5x^{2}+4 = 0$. En posant $X = x^{2}$ : $X^{2}-5X+4 = 0$, de
discriminant $9$, donc $X = 1$ ou $X = 4$. Les deux sont positives :
\[
  x \in \{-2\,;\,-1\,;\,1\,;\,2\}.
\]
\end{exemple}

\subsection{L'équation du troisième degré à racine évidente}

\begin{methode}{résoudre une équation de degré 3}
\begin{enumerate}[label=\textbf{\arabic*.}, leftmargin=*, itemsep=0.15em]
  \item Chercher une \textbf{racine évidente} parmi $-2$, $-1$, $0$, $1$,
        $2$, $3$ : on essaie, on ne devine pas.
  \item Si $x_{0}$ est racine, factoriser :
        $P(x) = (x-x_{0})\left(\alpha x^{2}+\beta x+\gamma\right)$. Les
        coefficients $\alpha$, $\beta$, $\gamma$ s'obtiennent par
        identification ou par division.
  \item Résoudre le trinôme obtenu.
\end{enumerate}
\end{methode}

\begin{exemple}[identification des coefficients]
Résolvons $x^{3}-2x^{2}-5x+6 = 0$. On essaie $x = 1$ :
$1-2-5+6 = 0$, donc $1$ est racine. On écrit alors
\[
  x^{3}-2x^{2}-5x+6 = (x-1)\left(x^{2}+\beta x-6\right),
\]
le coefficient de $x^{2}$ valant $1$ et le terme constant $-6$ (puisque
$(-1)\times(-6) = 6$). En développant, le coefficient de $x^{2}$ vaut
$\beta-1$, qui doit être égal à $-2$ : donc $\beta = -1$. Il reste
$x^{2}-x-6 = 0$, de racines $-2$ et $3$. L'ensemble des solutions est
$\{-2\,;\,1\,;\,3\}$.
\end{exemple}

\subsection{L'équation quotient}

\begin{avertissement}
Une équation quotient commence par la recherche des \textbf{valeurs
interdites} : celles qui annulent un dénominateur. Elles sont exclues avant
tout calcul, et une solution qui coïncide avec l'une d'elles doit être
rejetée à la fin.
\end{avertissement}

\begin{exemple}[valeurs interdites d'abord]
Résolvons $\dfrac{x+3}{x-1} = x$. La valeur $1$ est interdite. Pour
$x \neq 1$, l'équation équivaut à $x+3 = x(x-1)$, soit $x^{2}-2x-3 = 0$, de
racines $-1$ et $3$. Aucune n'est interdite :
$\mathcal S = \{-1\,;\,3\}$.
\end{exemple}

\begin{entrainement}[équations ramenées au second degré]
\exo[\niveau{1}] $x^{4}-13x^{2}+36 = 0$

\exo[\niveau{2}] $x^{4}+3x^{2}-4 = 0$

\exo[\niveau{2}] $x^{3}-x^{2}-4x+4 = 0$ (racine évidente : $1$)

\exo[\niveau{2}] $\dfrac{2x-1}{x+2} = 3$

\exo[\niveau{3}] $x^{3}-3x+2 = 0$

\exo[\niveau{3}] $\dfrac{1}{x}+\dfrac{1}{x+1} = \dfrac{5}{6}$
\end{entrainement}

\begin{corrige}[équations ramenées au second degré]
\textbf{1.} $X^{2}-13X+36 = 0$ donne $X = 4$ ou $X = 9$, d'où
$x \in \{-3\,;\,-2\,;\,2\,;\,3\}$.
\textbf{2.} $X^{2}+3X-4 = 0$ donne $X = 1$ ou $X = -4$ ; seule $X = 1$
convient, d'où $x = \pm1$.
\textbf{3.} $(x-1)(x^{2}-4) = 0$ : $x \in \{-2\,;\,1\,;\,2\}$.
\textbf{4.} Valeur interdite $-2$ ; $2x-1 = 3x+6$ donne $x = -7$.
\textbf{5.} $1$ est racine évidente ; $x^{3}-3x+2 = (x-1)(x^{2}+x-2)
= (x-1)^{2}(x+2)$, d'où $x \in \{-2\,;\,1\}$.
\textbf{6.} Valeurs interdites $0$ et $-1$. En réduisant :
$6(2x+1) = 5x(x+1)$, soit $5x^{2}-7x-6 = 0$, $\Delta = 169$,
$x = -\frac35$ ou $x = 2$.
\end{corrige}

\begin{vraifaux}
\exo Si $\Delta < 0$, le trinôme n'a pas de racine réelle, donc il change de
signe.

\exo Un trinôme dont le discriminant est nul garde un signe constant sur
$\R$.

\exo L'équation $x^{2} = x$ a pour unique solution $x = 1$.

\exo Si $x_{1}$ et $x_{2}$ sont les racines de $ax^{2}+bx+c$, alors
$x_{1}+x_{2} = \dfrac ba$.

\exo Une équation bicarrée a toujours quatre solutions.
\end{vraifaux}

\begin{corrige}[vrai ou faux]
\textbf{1.} Faux — c'est l'inverse : sans racine, il ne peut pas changer de
signe ; il garde celui de $a$.
\textbf{2.} Vrai — il s'écrit $a(x-x_{0})^{2}$, du signe de $a$, et s'annule
seulement en $x_{0}$.
\textbf{3.} Faux — $x^{2}-x = x(x-1) = 0$ donne $0$ et $1$. Diviser les deux
membres par $x$ fait perdre la solution $0$ : on ne divise jamais par une
quantité qui peut être nulle.
\textbf{4.} Faux — la somme vaut $-\frac ba$.
\textbf{5.} Faux — elle en a quatre, deux, une ou aucune, selon les signes
des solutions en $X$.
\end{corrige}

% ===========================================================================
\begin{situation}[le champ rectangulaire]
Une famille dispose de $86$ mètres de clôture pour entourer
un champ rectangulaire de $450$ mètres carrés, appuyé sur aucun mur : les
quatre côtés doivent être clôturés.

\begin{enumerate}[label=\textbf{\arabic*.}, leftmargin=*, itemsep=0.3em]
  \item On note $x$ la longueur et $y$ la largeur, en mètres. Traduire les
        deux contraintes par deux équations.
  \item Montrer que $x$ est solution de l'équation
        $x^{2}-43x+450 = 0$.
  \item Résoudre cette équation et donner les dimensions du champ.
  \item La famille voudrait, avec la même longueur de clôture, un champ de
        $500$ mètres carrés. Montrer que c'est impossible.
\end{enumerate}
\end{situation}

\begin{corrige}[le champ rectangulaire]
\textbf{1.} Le périmètre donne $2(x+y) = 86$, soit $x+y = 43$ ; l'aire donne
$xy = 450$.

\textbf{2.} $x$ et $y$ ont pour somme $43$ et pour produit $450$ : ce sont
les racines de $t^{2}-43t+450 = 0$.

\textbf{3.} $\Delta = 1849-1800 = 49$, $\sqrt\Delta = 7$ : les racines sont
$18$ et $25$. Le champ mesure $25$ m sur $18$ m.

\textbf{4.} Il faudrait $S = 43$ et $P = 500$, donc
$\Delta = 1849-2000 = -151 < 0$ : aucun couple de réels ne convient. Avec
$86$ mètres de clôture, l'aire maximale possible est celle du carré de côté
$21{,}5$ m, soit $462{,}25$ mètres carrés — et $500$ dépasse ce maximum.
\end{corrige}

% ===========================================================================
\begin{annales}
Le second degré n'est presque jamais l'objet d'un exercice entier à
l'examen : il apparaît \emph{à l'intérieur} des questions, au moment
d'étudier le signe d'une dérivée ou de résoudre une équation issue d'un
problème. Les énoncés ci-dessous en sont des extraits authentiques, réduits à
la partie qui ne dépend d'aucune notion postérieure à ce chapitre.

\exoann{Baccalauréat, Madagascar, série A, session 2026 — problème,
question 2.c) (extrait)}
Étudier le signe de $-2x^{2}+x+1$ pour tout $x$ de $\intervalleo{0}{+\infty}$.

\exoann{Baccalauréat, Madagascar, série A, session 2020 — problème,
question 2 (extrait)}
Étudier le signe du produit $(1+x)(3-x)$ sur l'intervalle
$\intervalleo{-1}{3}$, puis dire ce qu'il advient de ce signe en dehors de
cet intervalle.

\exoann{D'après le baccalauréat blanc, Lycée Philibert Tsiranana, série A,
2023 — problème, question 4.a)}
\emph{Le sujet pose cette question sur une expression construite au
chapitre 8 ; l'inconnue est notée ici $T$.}

Vérifier que $T^{2}-4T+3 = (T-1)(T-3)$, puis résoudre $T^{2}-4T+3 = 0$ et
donner le signe de $T^{2}-4T+3$ suivant les valeurs de $T$.
\end{annales}

\begin{corrige}[annales]
\textbf{1.} $\Delta = 1+8 = 9$ ; les racines sont
$\frac{-1-3}{-4} = 1$ et $\frac{-1+3}{-4} = -\frac12$. Comme $a = -2 < 0$, le
trinôme est positif entre les racines. Sur
$\intervalleo{0}{+\infty}$, il est donc strictement positif sur
$\intervalleo{0}{1}$, nul en $1$, strictement négatif sur
$\intervalleo{1}{+\infty}$.

\textbf{2.} Le produit s'annule en $-1$ et en $3$. Développé, il vaut
$-x^{2}+2x+3$, dont le coefficient dominant est négatif : il est donc positif
entre les racines et négatif à l'extérieur. Sur $\intervalleo{-1}{3}$, il est
strictement positif ; en dehors de $\intervalle{-1}{3}$, strictement négatif.

\textbf{3.} Le développement donne $T^{2}-3T-T+3 = T^{2}-4T+3$ : l'égalité
est vérifiée. Les solutions sont $T = 1$ et $T = 3$. Le coefficient dominant
étant positif, l'expression est positive à l'extérieur de
$\intervalle{1}{3}$, négative à l'intérieur.
\end{corrige}

\begin{remarque}[où retrouver le second degré]
Ce chapitre-ci ne peut pas encore recevoir les énoncés complets : ceux du
baccalauréat mêlent le trinôme au logarithme, à l'exponentielle ou à la
dérivée. Vous les retrouverez à leur place — au chapitre 5 pour le signe
d'une dérivée, au chapitre 6 pour l'étude d'une fonction rationnelle, aux
chapitres 7 et 8 pour les équations en $\ln x$ et en $e^{x}$, où le
changement d'inconnue ramène chaque fois à un trinôme.
\end{remarque}

% ===========================================================================
\begin{jecherche}[un paramètre dans l'équation]
On considère l'équation $(E_{m})$ : $x^{2}-2mx+m+2 = 0$, où $m$ est un réel.

\begin{enumerate}[label=\textbf{\arabic*.}, leftmargin=*, itemsep=0.3em]
  \item Calculer le discriminant de $(E_{m})$ en fonction de $m$.
  \item Pour quelles valeurs de $m$ l'équation admet-elle deux solutions
        distinctes ?
  \item Pour quelles valeurs de $m$ n'admet-elle aucune solution réelle ?
  \item Existe-t-il une valeur de $m$ pour laquelle $1$ est solution de
        $(E_{m})$ ? Si oui, déterminer alors l'autre solution.
\end{enumerate}
\end{jecherche}

\begin{corrige}[un paramètre dans l'équation]
\textbf{1.} $\Delta = 4m^{2}-4(m+2) = 4\left(m^{2}-m-2\right)$.

\textbf{2.} $\Delta > 0 \iff m^{2}-m-2 > 0$. Ce trinôme en $m$ a pour
discriminant $9$ et pour racines $-1$ et $2$ ; son coefficient dominant est
positif. Donc
$m \in \intervalleo{-\infty}{-1}\cup\intervalleo{2}{+\infty}$.

\textbf{3.} $\Delta < 0 \iff m \in \intervalleo{-1}{2}$.

\textbf{4.} $1$ est solution si $1-2m+m+2 = 0$, soit $m = 3$. L'équation
devient $x^{2}-6x+5 = 0$ ; le produit des racines vaut $5$, donc l'autre
racine est $5$.
\end{corrige}

% ===========================================================================
\begin{jemeteste}
Vingt minutes, sans le cours sous les yeux.

\begin{enumerate}[label=\textbf{\arabic*.}, leftmargin=*, itemsep=0.3em]
  \item Résoudre $2x^{2}-5x-3 = 0$.
  \item Factoriser $-x^{2}+4x-3$.
  \item Résoudre $x^{2}-x-12 \leq 0$.
  \item Résoudre $x^{4}-10x^{2}+9 = 0$.
  \item Déterminer les valeurs de $x$ pour lesquelles $\dfrac{x+1}{x^{2}-4}$
        est défini.
  \item Deux nombres ont pour somme $5$ et pour produit $-14$. Lesquels ?
\end{enumerate}
\end{jemeteste}

\begin{corrige}[je me teste]
\textbf{1.} $\Delta = 49$ ; $x = -\frac12$ ou $x = 3$.
\textbf{2.} Racines $1$ et $3$, $a = -1$ : $-(x-1)(x-3)$.
\textbf{3.} Racines $-3$ et $4$ : $\mathcal S = \intervalle{-3}{4}$.
\textbf{4.} $X = 1$ ou $X = 9$ : $x \in \{-3\,;\,-1\,;\,1\,;\,3\}$.
\textbf{5.} Défini pour $x \neq -2$ et $x \neq 2$.
\textbf{6.} Racines de $x^{2}-5x-14 = 0$ : $-2$ et $7$.
\end{corrige}

% ===========================================================================
\begin{commeaubac}
\bmtsource{Exercice original au format de l'épreuve — série A, options A1 et
A2 \textbullet\ 5 points}
Aucun sujet de la série n'a consacré d'exercice entier au second degré : le
sujet ci-dessous est composé au format et au barème de l'épreuve.

\begin{enumerate}[label=\textbf{\arabic*.}, leftmargin=*, itemsep=0.28em]
  \item On considère le polynôme $P(x) = x^{3}-4x^{2}+x+6$.
    \begin{enumerate}[label=\textbf{\alph*)}, leftmargin=1.4em, itemsep=0.15em]
      \item Vérifier que $P(-1) = 0$. \bareme{0,5 pt}
      \item Déterminer les réels $b$ et $c$ tels que, pour tout réel $x$,
            $P(x) = (x+1)\left(x^{2}+bx+c\right)$. \bareme{1,5 pt}
      \item Résoudre dans $\R$ l'équation $P(x) = 0$. \bareme{1 pt}
    \end{enumerate}
  \item Résoudre dans $\R$ l'inéquation $P(x) \geq 0$. \bareme{1,5 pt}
  \item En déduire l'ensemble des réels $x$ pour lesquels $\sqrt{P(x)}$ est
        défini. \bareme{0,5 pt}
\end{enumerate}
\end{commeaubac}

\begin{corrige}[comme au baccalauréat]
\textbf{1. a)} $P(-1) = -1-4-1+6 = 0$.

\textbf{b)} En développant $(x+1)(x^{2}+bx+c) = x^{3}+(b+1)x^{2}+(c+b)x+c$ et
en identifiant : $c = 6$, $b+1 = -4$ donc $b = -5$, et l'on vérifie
$c+b = 1$. Ainsi $P(x) = (x+1)\left(x^{2}-5x+6\right)$.

\textbf{c)} $x^{2}-5x+6 = 0$ a pour racines $2$ et $3$ ; d'où
$\mathcal S = \{-1\,;\,2\,;\,3\}$.

\textbf{2.} Tableau de signes du produit $(x+1)(x-2)(x-3)$ : le polynôme est
négatif sur $\intervalleo{-\infty}{-1}$, positif sur $\intervalle{-1}{2}$,
négatif sur $\intervalle{2}{3}$, positif sur
$\intervallefo{3}{+\infty}$. Donc
$P(x) \geq 0$ pour
$x \in \intervalle{-1}{2}\cup\intervallefo{3}{+\infty}$.

\textbf{3.} La racine carrée est définie exactement là où $P(x) \geq 0$ :
même ensemble qu'à la question 2.
\end{corrige}

\begin{notepourlaclasse}
Six heures suffisent si l'on renonce à tout redémontrer : la forme canonique
peut être admise et sa démonstration lue à la maison. Le temps gagné est à
mettre sur le tableau de signes, qui reste, chaque année, le point faible des
copies — bien avant le calcul du discriminant.

Le bloc d'automatismes est à refaire en début de chaque séance pendant deux
semaines, chronomètre en main : trois minutes, huit calculs. C'est le
meilleur investissement du chapitre.

Deux questions méritent d'être traitées au tableau plutôt que données à
chercher : l'exercice 6 des inéquations (quotient et valeur interdite), où
les élèves oublient presque tous que le dénominateur ne se contente pas de
changer de signe, et la question 4 du bloc « je cherche », qui demande de
raisonner sur un paramètre — un mouvement de pensée nouveau pour la plupart.
\end{notepourlaclasse}
